\documentclass[11pt]{article}
\usepackage[margin=1in]{geometry}
\usepackage{amsmath,amssymb}
\usepackage{graphicx}
\usepackage{booktabs}
\usepackage{array}
\usepackage{xcolor}
\usepackage{hyperref}
\usepackage{enumitem}
\usepackage{caption}
\usepackage{fancyhdr}
\usepackage{titlesec}

\hypersetup{colorlinks=true, linkcolor=blue, urlcolor=blue, citecolor=blue}
\titleformat{\section}{\large\bfseries}{\thesection}{0.6em}{}
\titleformat{\subsection}{\normalsize\bfseries}{\thesubsection}{0.6em}{}
\setlist[itemize]{leftmargin=1.4em}

\title{\vspace{-1.5em}AMCAS Assignment 1\\[0.3em]\large Memory Circuits and Systems: From Devices to Systems}
\author{Vyshnavi Tanikonda \, (2025900014)}
\date{}

\begin{document}
\maketitle
\vspace{-2em}

\section*{Setup \& Configuration Details}
All simulation source code, analysis scripts, raw data traces, and plotting configurations are fully reproducible and publicly available in the following GitHub repository: \url{https://github.com/VYSHNAVI2469/AMCAS--ASG1--Memory-sim}.

All simulations were conducted natively on an x86\_64 Linux environment (Ubuntu, WSL2). The tool configurations are as follows:
\begin{itemize}
    \item \textbf{ngspice}: Predictive Technology Model (PTM) 45\,nm bulk BSIM4 model card, $V_{DD,nom}=1.1$\,V.
    \item \textbf{CACTI 7.0}: 45\,nm ITRS High-Performance (\texttt{itrs-hp}) configuration at $T=350$\,K, UCA model.
    \item \textbf{NVSim}: 1T-1MTJ STT-MRAM cell at 45\,nm, current-sensing read, current-mode set/reset.
    \item \textbf{Ramulator}: DDR4-2400 8x8, single channel, FR-FCFS scheduler.
    \item \textbf{gem5 (v24.0.0.0)}: x86 ISA, \texttt{DerivO3CPU} (Out-of-Order) and \texttt{TimingSimpleCPU} (In-Order) with the GAPBS benchmark suite.
\end{itemize}

\textit{Central architectural question carried through all five parts: \textbf{Should the 2\,MB L2 cache in our accelerator use SRAM or STT-MRAM?}}

\section{Part A: 6T SRAM Bitcell Physics (ngspice)}
The baseline 6T SRAM bitcell uses $V_{DD}=1.1$\,V. The pull-down (MN1) to access-transistor (MA1) sizing ratio is $\beta_r = W_{pd}/W_{ax} = 0.20/0.16 = 1.25$. The cell is initialized with $Q=0$\,V, $QB=1.1$\,V, and both bitlines precharged to $1.1$\,V with $C_{BL}=180$\,fF.

\textbf{Task 1 (Baseline read):} At $t=2.0$\,ns, the measured differential is $\Delta V(BLB,BL) = \mathbf{729.86\ mV}$, with the internal storage node peaking at $v(q)_{max}=\mathbf{205.52\ mV}$ at $t=1.058$\,ns. This is roughly $10.6\times$ the $\approx69$\,mV signal a 1T1C DRAM cell develops from passive charge sharing alone, since the SRAM cell \emph{actively} discharges the bitline through MA1/MN2 rather than relying on charge redistribution.

\textbf{Task 2 (Widened access transistor):} Growing the access transistors to $W_{ax}=0.24\,\mu$m drops the cell ratio to $\beta_r = 0.20/0.24 = 0.833 < 1.0$. The storage-node bump grows to $v(q)_{max}\approx\mathbf{272.0\ mV}$, consuming more than 60\% of the inverter trip voltage ($V_{trip}\approx450$\,mV). Under the ideal, noise-free transient conditions simulated, the disturbance still relaxes back toward $0$\,V rather than latching into a destructive flip -- but with realistic threshold-mismatch ($\sigma_{V_{th}}\approx35$\,mV from RDF), cells in the distribution tail would very plausibly trip \texttt{MN1} and undergo a destructive read upset. (See Figure~\ref{fig:parta_transient}.)

\begin{figure}[h]
    \centering
    \includegraphics[width=0.95\textwidth]{figs/fig_partA_transient.png}
    \caption{Bitline differential and storage-node transients for $W_{ax}=0.16\,\mu$m (baseline, $\beta_r=1.25$) vs. $W_{ax}=0.24\,\mu$m (widened, $\beta_r=0.833$) (Task 1/2).}
    \label{fig:parta_transient}
\end{figure}

\textbf{Task 3 ($V_{DD}$ sweep, $C_{BL}=1$\,pF):} Sweeping $V_{DD}$ from $1.1$\,V down to $0.6$\,V, $\Delta V(BLB,BL)$ at $t=2.0$\,ns falls monotonically from $143.59$\,mV to $26.91$\,mV. Using Lecture 5's $25$\,mV sense-amplifier offset and interpolating between the $0.65$\,V and $0.60$\,V points gives a failure boundary of $V_{DD}\approx\mathbf{0.591\ V}$ -- below this, the developed bitline differential can no longer be reliably resolved by the sense amplifier. (See Figure~\ref{fig:parta_vdd}.)

\begin{figure}[h]
    \centering
    \includegraphics[width=0.72\textwidth]{figs/fig_partA_vdd.png}
    \caption{Bitline differential $\Delta V$ vs. supply voltage $V_{DD}$, $C_{BL}=1$\,pF (Task 3).}
    \label{fig:parta_vdd}
\end{figure}

\textbf{Task 4 (85\,\textdegree C corner):} At $85\,^{\circ}$C, $\Delta V(BLB,BL)$ at $2.0$\,ns falls from $729.86$\,mV to $\mathbf{582.11\ mV}$ (a $20.2\%$ drop) due to reduced carrier mobility ($\mu\propto T^{-1.5}$), while the internal read bump \emph{increases} slightly, from $205.52$\,mV to $220.76$\,mV ($+7.4\%$). Because the raw bitline swing still comfortably clears the $25$\,mV sense-amp offset at this corner, the dominant failure mechanism at high temperature is not the swing itself but the combination of temperature-dependent sense-amplifier offset drift and rising subthreshold leakage from unselected cells sharing the bitline.

\section{Part B: SRAM Geometry \& Capacity Sizing (CACTI 7)}

\textbf{Task 1 (2\,MB baseline, ED\textsuperscript{2}P):}
\begin{center}
\begin{tabular}{lr}
\toprule
Metric & Value \\
\midrule
Access time & $\mathbf{2.9018}$ ns \\
Cycle time & $2.6523$ ns \\
Dynamic read energy & $792.89$ pJ \\
Dynamic write energy & $851.22$ pJ \\
Leakage power (per bank / total 4 banks) & $562.63$ mW / $\mathbf{2250.53}$ mW \\
Total area & $\mathbf{11.474}$ mm$^2$ ($2.213\times5.184$\,mm) \\
Winning organization $(N_{dwl},N_{dbl},N_{spd})$ & $(4,2,1)$ \\
\bottomrule
\end{tabular}
\end{center}
Of the $2.9018$\,ns access time, $46.7\%$ ($1.36$\,ns) is consumed by global H-tree address/data routing alone -- more than the decoder, bitline, and sense-amp stages combined.

\textbf{Task 2 (Capacity sweep, $256$\,kB $\to$ $16$\,MB):} Access time nearly triples ($2.417\to6.362$\,ns) across a $64\times$ capacity range (Figure~\ref{fig:partb_sweep}). Between $256$\,kB and $512$\,kB the increment ($+67.4$\,ps) tracks Amrutur \& Horowitz's ``one gate delay per doubling'' rule; beyond $\approx2$\,MB the per-doubling delay accelerates superlinearly ($+269$\,ps $\to$ $+1890$\,ps) because global H-tree wire length scales as $\sqrt{\text{Area}}$ and its RC delay overwhelms the constant gate-delay term.

\begin{figure}[h]
    \centering
    \includegraphics[width=0.85\textwidth]{figs/fig_partB_sweep.png}
    \caption{CACTI access latency vs. capacity (log$_2$ scale), $256$\,kB$-16$\,MB (Task 2).}
    \label{fig:partb_sweep}
\end{figure}

\textbf{Task 3 (Objective function comparison):}
\begin{center}
\begin{tabular}{lccc}
\toprule
Objective & $(N_{dwl},N_{dbl},N_{spd})$ & Access Time & Full-Cache Area \\
\midrule
Pure Delay ($100{:}0{:}0{:}0{:}0$) & $(4,2,1)$ & $2.892$ ns & $10.970$ mm$^2$ \\
Pure Area ($0{:}0{:}0{:}0{:}100$) & $(2,2,1)$ & $3.355$ ns & $\mathbf{9.245}$ mm$^2$ ($-19.4\%$) \\
ED\textsuperscript{2}P (baseline) & $(4,2,1)$ & $2.902$ ns & $11.474$ mm$^2$ \\
\bottomrule
\end{tabular}
\end{center}
Pure Delay and ED\textsuperscript{2}P agree on $(4,2,1)$ because ED\textsuperscript{2}P's $\text{Delay}^2$ term heavily penalizes latency; Pure Area instead merges wordline partitions to $N_{dwl}=2$, sharing row decoders at the cost of longer wordlines and $+16\%$ delay.

\textbf{Task 4 (Hand-check vs. CACTI bitline delay):} The Elmore-model hand calculation $\tau=0.38RC\approx1.16$\,ps is roughly $350\times$ smaller than CACTI's reported $406.9$\,ps. The gap is explained by effects the lumped hand model ignores: the access-transistor's ON-resistance ($8$--$12$\,k$\Omega$, $\gg$ metal wire resistance), junction capacitance from $\approx512$ unselected cells on the column, mux/sense-amp/precharge loading, and CACTI's use of a small-signal sensing threshold ($\Delta V\approx100$\,mV) rather than a rail-to-rail $50\%$ step.

\section{Part C: STT-MRAM Evaluation (NVSim)}

\textbf{Task 1 (2\,MB STT-MRAM vs. SRAM baseline):} MTJ cell: $54\,F^2$ area, aspect ratio $2.0$, $R_{on}=3$\,k$\Omega$ / $R_{off}=6$\,k$\Omega$ (TMR $2{:}1$), read current $40\,\mu$A, set/reset current $200\,\mu$A over a $10$\,ns pulse.

\begin{center}
\begin{tabular}{lrrr}
\toprule
Metric & SRAM (CACTI) & STT-MRAM (NVSim) & Ratio \\
\midrule
Read (hit) latency & $2.902$ ns & $\mathbf{1.589}$ ns & $0.55\times$ \\
Write latency & $2.652$ ns$^{*}$ & $\mathbf{10.608}$ ns & $4.00\times$ \\
Read dynamic energy & $0.793$ nJ & $0.760$ nJ & $0.96\times$ \\
Write dynamic energy & $0.851$ nJ & $0.531$ nJ & $0.62\times$ \\
Leakage power & $2250.5$ mW & $\mathbf{425.2}$ mW & $0.19\times$ \\
Total area & $11.474$ mm$^2$ & $\mathbf{3.590}$ mm$^2$ & $0.31\times$ \\
\bottomrule
\end{tabular}
\end{center}
\footnotesize{$^{*}$SRAM ``write latency'' approximated by CACTI's reported cycle time; CACTI does not separately report an asymmetric write path.}
\normalsize

\begin{figure}[h]
    \centering
    \includegraphics[width=0.85\textwidth]{figs/fig_partC_compare.png}
    \caption{Normalized comparison of 2\,MB STT-MRAM vs. 2\,MB SRAM baseline (Task 1).}
    \label{fig:partc_compare}
\end{figure}

\textbf{Task 2 (What flips, and why):} \emph{Dramatically better:} area ($0.31\times$) and leakage ($0.19\times$) -- the 1T-1MTJ cell needs one access transistor instead of six, and the MTJ holds state magnetically, so standby current is almost entirely peripheral CMOS rather than the storage element. \emph{Dramatically worse:} write latency ($4.0\times$) -- switching the MTJ's free-layer magnetization requires sustaining a $200\,\mu$A spin-transfer-torque current for a full $10$\,ns pulse, versus a single cycle to overpower a bistable 6T latch.

\textbf{Task 3 (TMR re-run, $R_{on}=4$\,k$\Omega$ / $R_{off}=12$\,k$\Omega$, TMR $3{:}1$):} Nearly every array-level metric is essentially unchanged (area $3.590\to3.580$\,mm$^2$; hit latency $1.589\to1.594$\,ns; write latency $10.608\to10.657$\,ns; leakage unchanged at $425.227$\,mW). Because NVSim uses a fixed \texttt{MinSenseVoltage} and current-mode sensing, the $2{:}1$ TMR already satisfies the sense margin comfortably -- a higher TMR buys additional \emph{read-margin reliability} headroom, not fundamentally faster or smaller arrays at this abstraction level.

\textbf{Task 4 (Halving ResetCurrent to $100\,\mu$A):} Write current is the key coupling variable for NVM cell area: a lower reset/set current permits a smaller access transistor, which in principle shrinks the 1T-1MTJ cell. Because NVSim treats cell area as a fixed parameter from the \texttt{.cell} file rather than deriving it from the specified currents, the reported area does not update automatically -- the designer must manually recompute and pass in the smaller \texttt{-CellArea} to see the resulting density benefit. This is an important asymmetry vs. SRAM, where cell area is fixed by the 6T ratioing constraints rather than by a write-current specification.

\section{Part D: DRAM Memory Scheduling (Ramulator)}
Ramulator replays the L2-miss address stream from the gem5/CACTI flow above against a single-channel DDR4-2400 8x8 DIMM with an FR-FCFS scheduler, open-loop (no processor-stall feedback).

\begin{itemize}
    \item \textbf{Scheduler-bound behavior:} the channel is bandwidth/$t_{CCD}$-limited -- i.e., bounded by the data-bus burst cycle time rather than by row-activation ($t_{RCD}/t_{RAS}$) timing -- so \emph{reducing the number of L2 misses} (the job Part C's larger STT-MRAM capacity does) matters more for end-to-end throughput than shaving individual miss latency.
    \item \textbf{Queue occupancy:} the open-loop replay's request queue sits around $36$--$50$ requests, reflecting the fact that misses are injected without waiting on ROB/processor backpressure.
    \item \textbf{Methodological caveat:} because the replay is open-loop, it can overstate queue occupancy and understate how much latency the real out-of-order core would otherwise hide -- this is one of the three assumptions flagged for review in Part E, Task 4.
\end{itemize}

\section{Part E: Full-System Execution (gem5) \& Synthesis}
Setup: \texttt{DerivO3CPU} (OoO, 2.0\,GHz, 8-wide issue, 192-entry ROB) and \texttt{TimingSimpleCPU} (in-order); 32\,kB, 8-way L1 I/D caches (2-cycle hit); DDR4-2400 main memory; GAPBS \texttt{bfs}/\texttt{sssp} on a scale-18 Kronecker graph (262{,}143 vertices, 3{,}805{,}449 edges). The 2\,MB SRAM L2 (7-cycle hit, from Part B) is compared against an area-matched 8\,MB STT-MRAM L2 (14-cycle hit, $\approx11.11$\,mm$^2$ from Part C, vs. $11.47$\,mm$^2$ for the SRAM baseline).

\textbf{Task 1 (Out-of-order results):}
\begin{center}
\begin{tabular}{llrrrr}
\toprule
Kernel & L2 Config & IPC & L2 Miss Rate & simSeconds & Speedup \\
\midrule
bfs & SRAM 2\,MB @ 7cy & $0.5266$ & $68.7\%$ & $0.015496$ & $1.000\times$ \\
bfs & MRAM 8\,MB @ 14cy & $0.7420$ & $18.2\%$ & $0.010998$ & $\mathbf{1.409\times}$ \\
sssp & SRAM 2\,MB @ 7cy & $0.4178$ & $60.5\%$ & $0.110702$ & $1.000\times$ \\
sssp & MRAM 8\,MB @ 14cy & $0.4469$ & $44.7\%$ & $0.103506$ & $\mathbf{1.070\times}$ \\
\bottomrule
\end{tabular}
\end{center}

\textbf{Task 2 (Why BFS wins big, SSSP barely):} BFS's hot working set (parent array + visited bitmap + frontier buffers, $\approx2.2$--$3.5$\,MB) just barely overflows 2\,MB but fits comfortably in 8\,MB, so the miss rate drops $3.77\times$ ($68.7\%\to18.2\%$) and each avoided DRAM miss ($\approx140$--$160$ cycles) vastly outweighs the extra 7-cycle hit penalty. SSSP's $\Delta$-stepping access pattern is far more scattered ($7.7\times$ more L2 accesses than BFS) and its footprint does not fit meaningfully better in 8\,MB, so the miss-rate improvement is smaller ($1.35\times$) while over $1.8$M L2 \emph{hits} now each pay the extra latency -- largely cancelling the DRAM-miss savings.

\textbf{Task 3 (In-order core):}
\begin{center}
\begin{tabular}{llrrrr}
\toprule
Kernel & L2 Config & IPC & Miss Rate & simSeconds & Speedup \\
\midrule
bfs & SRAM & $0.1803$ & $69.0\%$ & $0.045263$ & $1.000\times$ \\
bfs & MRAM & $0.2401$ & $15.9\%$ & $0.033986$ & $1.332\times$ \\
sssp & SRAM & $0.1505$ & $60.8\%$ & $0.307957$ & $1.000\times$ \\
sssp & MRAM & $0.1601$ & $45.1\%$ & $0.289435$ & $1.064\times$ \\
\bottomrule
\end{tabular}
\end{center}
The 192-entry ROB in \texttt{DerivO3CPU} overlaps much of the extra 7-cycle STT-MRAM hit latency with independent work, so BFS's speedup is \emph{larger} under OoO ($1.409\times$) than in-order ($1.332\times$) -- OoO hides \emph{hit} latency but cannot hide the $\approx140$--$180$ cycle DRAM-\emph{miss} latency, which is why avoiding misses stays valuable on either core.

\begin{figure}[h]
    \centering
    \includegraphics[width=0.85\textwidth]{figs/fig_partE_speedup.png}
    \caption{Full-system speedup of 8\,MB STT-MRAM over 2\,MB SRAM, OoO vs. in-order cores.}
    \label{fig:parte_speedup}
\end{figure}

\subsection*{Architectural Synthesis \& Deliverables (Task 4)}

\textbf{Central question answer:} \emph{Should the 2\,MB L2 cache in our accelerator be SRAM, or STT-MRAM?} Based on the cross-layer evaluation, the area-matched 8\,MB STT-MRAM L2 is $\mathbf{1.409\times}$ faster than the 2\,MB SRAM baseline on the BFS workload with an out-of-order core. While STT-MRAM suffers a $4\times$ worse write latency (NVSim, Part C) and a $2\times$ nominal read-hit latency (14 vs. 7 cycles), its $\approx3.2\times$ area density advantage lets the 8\,MB capacity collapse the L2 miss rate from $68.7\%$ to $18.2\%$ (gem5, Part E), which keeps the DDR4 channel from being swamped by miss traffic that Part D shows is already $t_{CCD}$-bound -- provided the core has enough out-of-order resources to mask the slower array response.

\textbf{Three assumptions a reviewer should challenge:}
\begin{enumerate}[leftmargin=1.6em]
    \item \textbf{The fixed 2$\times$ hit-latency assumption (14 vs. 7 cycles).} CACTI's $2.902$\,ns SRAM access and NVSim's $1.589$\,ns STT-MRAM read latency (at matched 2\,MB capacity) do not actually support a clean $2\times$ ratio when modeled consistently -- the $14$-cycle figure used in gem5 is a pedagogical simplification, not a directly derived physical result.
    \item \textbf{Open-loop Ramulator replay has no processor feedback} (Part D). Requests are generated without backpressure from stalls or the ROB, which can distort queue occupancy ($36$--$50$ requests) and over/understate how much of the real memory latency an actual core would hide.
    \item \textbf{Write endurance, wear-leveling, and thermal sensitivity are unmodeled.} Graph analytics such as SSSP generate intense, localized write traffic that NVSim/gem5 do not penalize for finite STT-MRAM write endurance; and Part A shows SRAM margins already erode at $85\,^{\circ}$C, while STT-MRAM retention is exponentially temperature-sensitive, so an 8\,MB high-performance L2 could face thermal-driven retention failure or refresh overhead the models omit. Finally, all of the above rests on a single synthetic Kronecker graph topology -- real-world sparse graphs (road networks, web graphs) have different degree distributions and working-set behavior, so the STT-MRAM recommendation should not be generalized beyond BFS-like access patterns without further validation.
\end{enumerate}

\end{document}
